\section{Conclusion}
\label{sec:conclusion}

\subsection{Summary}

We conducted a 50-state parameter sweep of five key \ar{} + \cs{}
parameters, finding that:

\begin{enumerate}
\item \textbf{Partisan outcomes are parameter-insensitive.}
      The maximum national partisan variation across the full parameter
      range is $\Delta D_{\rm nat} \leq 0.3$ seats---$60\times$ smaller
      than the algorithm-structure effect from B.0.
\item \textbf{Compactness is sensitive to $\ufactor$.}
      The METIS imbalance factor drives $\overline{\PP}$ variation of
      $\pm 0.010$ (3.2\% relative change).
      The DIA default $\ufactor = 1\%$ is the recommended statutory value.
\item \textbf{$T$ affects only runtime.}
      For $T \geq 200$, the \cs{} threshold has no effect on outcomes.
      The DIA default $T = 600$ provides a certified 89-seed safety margin
      (B.16) with negligible runtime cost.
\item \textbf{County weight affects county integrity.}
      $\acounty$ drives county split counts but has $\leq 0.2$-seat
      partisan effect nationally.
\end{enumerate}

\subsection{Policy Recommendation}

The DIA can safely encode the five parameters in statute at the
recommended defaults (Table~\ref{tab:recommended-defaults}).
No future legislature or administrator can achieve a meaningful partisan
outcome by lobbying for parameter changes.
The legitimate lever for partisan outcome---algorithm structure selection---was
decided at the statutory drafting stage (B.0) and is not subject to
administrative discretion.

\subsection{Limitations}

\begin{enumerate}
\item The OAT sweep does not capture parameter interactions.
      A full factorial design ($5^5 = 3{,}125$ runs) would be more
      complete but computationally expensive.
      Additive bound suggests interactions cannot change the conclusion.
\item The partisan metric uses 2020 election data.
      Future election outcomes may change which districts are competitive,
      potentially affecting the partisan sensitivity of some parameters.
\item We study only the parameters listed in Section~\ref{sec:parameters}.
      Other implementation choices (e.g., the spanning-tree sampler in
      \recom{}, the METIS initial partitioning method) could in principle
      be tuned for partisan effect.
\end{enumerate}

\subsection{Future Work}

A full factorial sweep of all five parameters would confirm the additive
bound and potentially identify interaction effects.
We leave this as future work, noting that the additive bound of 0.5 seats
is already small enough to be policy-irrelevant.
